\section{Conclusion}

This project implemented and evaluated value-prediction mechanisms
in gem5 MinorCPU, including Last Value Prediction (LVP),
ASTRA-based load-base prediction, a dedicated value cache, and
guarded V-ASTRA load bypass. It also investigated two
dependency-handling mechanisms: Design~A releases scoreboard
dependencies after a successful bypass, while Design~B forwards
validated values to eligible dependent loads.

The paper's META predictor was adapted to MinorCPU to select
between LVP and V-ASTRA on a per-consumer basis. Validated
V-ASTRA values are preferred, while consumers wait when a
V-ASTRA result is Pending rather than immediately falling back
to LVP. LVP may be selected when the V-ASTRA outcome has failed
or is absent and the LVP admission conditions are satisfied. The implementation also introduces experimental
provisional-base modes that use LVP-derived information for
ASTRA training and V-ASTRA candidate generation.

Performance measurements on deterministic microbenchmarks show
that prediction effectiveness depends on workload characteristics.
META achieves a 31.20\% cycle reduction on mixed,
outperforming both standalone predictors. However, on
random base, default META achieves 30.27\%, compared with
48.75\% for V-ASTRA-only, and it also performs slightly worse
than V-ASTRA on stable flip. The experimental candidate
mode improves the random-base reduction to 50.80\%, but this
benefit is workload-specific and does not generalize across the
evaluated workloads. These results characterize the tested simulator
configurations, not application-level speedups. Design~A was
validated under stall-on-use but not performance-evaluated in
that configuration, while the contribution of Design~B was not
isolated from V-ASTRA.

Validation covered predictor behavior, value-cache correctness,
memory ordering, dependency handling, and misprediction recovery.
The experiments demonstrated that unused incorrect LVP predictions
do not require recovery and that V-ASTRA-served consumers do not
trigger LVP recovery. They also exposed and helped correct a
provisional-state cleanup defect on early alignment faults.

Overall, this work demonstrates the challenges of integrating
value prediction into an in-order processor, particularly the
interactions between speculative operands, address prediction,
value-cache bypass, memory ordering, dependency tracking, and
recovery. It provides a basis for further investigation of value
prediction in MinorCPU. Important next steps include completing
final-source validation, implementing architectural components
such as the speculative L0 minicache and LACT, improving provisional
candidate admission, and evaluating the mechanisms on more
representative workloads.
